% !TEX root = ../../main.tex
\chapter{写在前面的话}

欢迎来到自动控制理论的殿堂！这是自动化专业最重要的专业课之一，虽然说学的内容大部分在后续的生产生活当中都用不到，有些过时了，但毕竟这是前人智慧的结晶！还是很重要的！

本笔记的主要参考书目为哈工大出版社出版的自动控制原理（上）和（下），参考的上课ppt为张宏伟老师、张颖老师、吴爱国老师和许鋆老师的ppt，他们分别教的内容为：
\begin{enumerate}
    \item 张宏伟老师：控制系统的运动分析（时域、频域，离散系统等等）
    \item 张颖老师：控制系统的稳定性分析
    \item 吴爱国老师：控制系统性能指标、频域校正、根轨迹校正
    \item 许鋆老师：PID、离散系统校正、能控能观、状态反馈以及状态观测、非线性控制分析
\end{enumerate}

不过笔者写笔记的思路和这些都有些不一样，本笔记主要分为五个部分：
\begin{enumerate}
    \item 经典控制理论（分析）
    \item 现代控制理论（分析）
    \item 系统的稳定性
    \item 控制系统设计综合
    \item 非线性系统的分析
\end{enumerate}

23级自动化用了三个学期来学习自动控制理论，不过好像后面的又改回了22级及以前的上课方式，采用自控理论A和自控理论B上课，这都无关紧要，学的内容都是大差不差的。祝愿各位哈工深的大三自动化学子们能度过一个愉快的大三生活！

也希望看到这份笔记的友友们如果发现有什么写错或者写漏了的部分可以向笔者提出来，笔者的联系方式为：\textbf{QQ：2745813833}，笔者不保证所有的内容都是对的，毕竟是上了三个学期才写出来的笔记，大家一起把这份笔记做完整做全对！

\begin{marker}
    本笔记中存在超链接，带有超链接的字会标粉，可以直接双击跳转到对应的位置
\end{marker}
